\documentclass[11pt,a4paper]{report}
\usepackage[a4paper,margin=24mm,headheight=14pt]{geometry}
\usepackage[T1]{fontenc}
\usepackage[utf8]{inputenc}
\usepackage{lmodern}
\usepackage{graphicx}
\usepackage{xcolor}
\usepackage{booktabs}
\usepackage{array}
\usepackage{tabularx}
\usepackage{longtable}
\usepackage{float}
\usepackage{caption}
\usepackage{fancyhdr}
\usepackage[hidelinks]{hyperref}
\usepackage{enumitem}
\usepackage{url}

\definecolor{ink}{RGB}{32,44,58}
\definecolor{accent}{RGB}{37,96,151}
\definecolor{soft}{RGB}{237,243,248}
\definecolor{warning}{RGB}{139,69,19}
\setlength{\parindent}{0pt}
\setlength{\parskip}{6pt}
\setlist[itemize]{leftmargin=18pt,itemsep=2pt,topsep=3pt}
\renewcommand{\arraystretch}{1.18}
\captionsetup{font=small,labelfont=bf}
\pagestyle{fancy}
\fancyhf{}
\lhead{\textcolor{ink}{Vision-Language Model Comparison}}
\rhead{\textcolor{ink}{\thepage}}
\renewcommand{\headrulewidth}{0.3pt}

\newcommand{\code}[1]{\texttt{#1}}
\newcommand{\figpath}[1]{comparison/output/#1}
\newcommand{\archbox}[2]{\fcolorbox{accent}{soft}{\parbox{#1}{\centering #2}}}
\newcommand{\resultbox}[1]{%
  \vspace{3pt}\noindent
  \fcolorbox{accent}{soft}{\parbox{0.94\textwidth}{#1}}\vspace{4pt}}
\newcommand{\twographs}[6]{%
\begin{figure}[H]
\centering
\begin{minipage}[t]{0.48\textwidth}
\centering
\includegraphics[width=\linewidth]{\figpath{#1}}
\captionof{subfigure}{#2}
\end{minipage}\hfill
\begin{minipage}[t]{0.48\textwidth}
\centering
\includegraphics[width=\linewidth]{\figpath{#3}}
\captionof{subfigure}{#4}
\end{minipage}
\caption{#5}
\label{#6}
\end{figure}}

\title{\textbf{Seeing the Same Task Differently}\\[8pt]
Vision-Language Model Operation, Benchmark Comparison,\\
and a Qwen2.5-VL Fine-Tuning Case Study}
\author{Technical report from the repository evaluation outputs}
\date{}

\begin{document}
\pagenumbering{roman}

\begin{titlepage}
\centering
\vspace*{34mm}
{\Huge\bfseries Seeing the Same Task Differently\par}
\vspace{10mm}
{\Large Vision-Language Model Operation, Benchmark Comparison,\\
and a Qwen2.5-VL Fine-Tuning Case Study\par}
\vspace{22mm}
\rule{0.68\textwidth}{0.5pt}\par
\vspace{12mm}
{\large A report built from the benchmark and fine-tuning artifacts\\
stored in this repository\par}
\vfill
{\normalsize Primary analysis command:\par}
\vspace{2mm}
\code{python comparison\textbackslash build\_report.py --results-dir results --output-dir comparison\textbackslash output --bootstrap-samples 5000}
\end{titlepage}
\setcounter{page}{2}

\chapter*{Abstract}
\addcontentsline{toc}{chapter}{Abstract}
This report starts from a practical question: if the same harness gives the
same kinds of images and prompts to several vision-language models, what
actually changes in their answers?  The setup in this repository is a good
place to ask that question.  Whether a model is loaded locally through Hugging
Face or called through an API wrapper, the outside contract is simple: an image
and a prompt go in, text comes out, and the benchmark turns that text into a
score plus some runtime information.

The main comparison covers 294 model--benchmark cells: fourteen models, twenty-
one benchmarks, and four task families: captioning, visual question answering
(QA), labeling, and detection.  On the benchmark-normalized summary,
\code{gpt-5.5} comes first at 0.825, with \code{gpt-4.1} second at 0.757.
That is the headline result, but it is not enough by itself.
\code{qwen2.5-vl-3b-instruct} has the most first-place finishes and is the
strongest captioning model in this run.  PaliGemma looks weak in the overall
mean, yet leads the normalized QA summary.  In practice, the right model
depends heavily on the kind of answer the application needs.

The report also includes a small Fashion-MNIST fine-tuning case study.
Balanced QLoRA adaptation of Qwen2.5-VL-3B improves held-out accuracy from
19/50 to 24/50.  That is a real gain on this sample, but it is not a clean win:
the adapted model predicts \emph{shirt} much more often and loses ground on
nearby clothing classes.  The useful lesson is narrow but important.  Fine-
tuning can move a small model, but the error pattern has to be checked just as
carefully as the headline accuracy.

\resultbox{\textbf{How to read this report.} Performance claims in the main
comparison come from \code{comparison/output/}, regenerated from the top-level
\code{results/} JSON files.  Fine-tuning results are handled separately; they
do not enter the aggregate leaderboard.}

\clearpage
\tableofcontents
\clearpage
\pagenumbering{arabic}

\chapter{Research Questions and Scope}

\section{Why compare through one harness?}
Vision-language models are usually introduced through parameter counts,
architecture names, or public benchmark scores.  Those are useful facts, but
they do not answer the question that matters inside this repository: given a
PIL image, a prompt, and an expected output format, which wrapper behaves best
when it is actually used in this benchmark code?

That kind of question needs a fairly ordinary kind of care.  The prompts have
to ask for a parseable answer.  Predictions have to be normalized in the same
way.  Captioning scores and classification accuracies cannot simply be averaged
as if they were the same unit.  Runtime numbers need to be kept, but they also
need context, because a hosted API call and a local GPU run are not the same
experiment.

The repository gives us a clean comparison point through a small shared
interface.  Every model wrapper exposes \code{predict(image, prompt)}: image
and prompt in, text out.  Every benchmark builds a prompt, sends it to the
wrapper, parses the answer, and records a score.  That is the level this report
trusts.  It compares what the wrappers do in this harness, not what a model
name may suggest in isolation.

\section{Questions addressed}
I organize the analysis around three questions.

\begin{enumerate}[leftmargin=22pt]
\item Which of the evaluated models is strongest overall, after respecting
      differences between benchmark metrics?
\item Do particular models have credible task specialisms, even when they do
      not top the aggregate table?
\item In a small, held-out clothing classification experiment, does QLoRA
      adaptation improve Qwen2.5-VL-3B, and what new mistakes appear?
\end{enumerate}

The first two questions use the main comparison output: fourteen model labels,
twenty-one benchmarks, and four task families.  The fine-tuning question uses
separate nested result files.  Once a model has been adapted, it is no longer
part of the same baseline ranking, so those results are kept out of the general
leaderboard.

\section{Evidence boundary}
The model labels in the figures are the result identifiers used by this
repository.  For local models, we can discuss loading, processing, and decoding
because the wrapper code is present.  For API-served models, the evidence is
more limited: the wrapper submits text plus a data-URL image and receives text
back.  The report therefore does not infer private architecture from an API
model name.

\begin{table}[H]
\centering
\caption{Evidence layers used in the report.}
\begin{tabularx}{\textwidth}{>{\bfseries}p{32mm} X X}
\toprule
Layer & Evidence & What it supports \\
\midrule
Implementation & \code{models/}, \code{benchmarks/}, \code{fine\_tuning/} &
Prompt, input, generation, parsing, and adaptation workflow statements. \\
Main results & CSV tables and PNG plots in \code{comparison/output/} &
Aggregate, family, benchmark, pairwise, uncertainty, and telemetry claims. \\
Adapter case study & JSON reports in \code{results/fine-tuning/} and the
downloaded smoke artifact folder & Paired Fashion-MNIST outcomes only. \\
\bottomrule
\end{tabularx}
\end{table}

\clearpage
\chapter{Models}

\section{One contract, several paths}
From the benchmark's point of view, every model looks similar: it receives an
image and a prompt, then returns text.  Inside the wrappers, the paths are very
different.  This chapter is about those visible paths in the repository.  It is
not trying to reconstruct private model internals or treat a model family name
as a complete architectural description.

\section{Model families, transformer structure, and paths}
There are two useful identifiers for each model: the wrapper path we can read,
and the checkpoint or API model ID that the wrapper calls.  For hosted API
models, the wrapper is the honest boundary of the evidence.  For local
checkpoints, the common pattern is easier to describe: visual information is
turned into embeddings or image tokens, then passed into a language model that
generates the answer.

\begin{table}[H]
\centering
\caption{Evaluated model labels covered by the architecture discussion.}
\small
\begin{tabularx}{\textwidth}{>{\bfseries}p{30mm} X}
\toprule
Architecture family & Model labels in the regenerated comparison \\
\midrule
OpenAI API vision pathway &
\code{gpt-4.1}, \code{gpt-4o}, \code{gpt-5}, \code{gpt-5.1},
\code{gpt-5.2}, \code{gpt-5.3-chat-latest}, \code{gpt-5.4},
\code{gpt-5.4-mini}, \code{gpt-5.4-nano}, \code{gpt-5.5}. \\
Qwen2.5-VL &
\code{qwen2.5-vl-3b-instruct}. \\
LLaVA-style local models &
\code{llava-1.5-7b-hf}, \code{llava-gemma-2b}. \\
PaliGemma &
\code{paligemma-3b-mix-224}. \\
\bottomrule
\end{tabularx}
\normalsize
\end{table}

\begin{table}[H]
\centering
\caption{Model paths and architecture level used in the report.}
\footnotesize
\begin{tabularx}{\textwidth}{>{\bfseries}p{25mm} p{38mm} p{39mm} X}
\toprule
Family & Repository path & Model ID or checkpoint path & Architecture described here \\
\midrule
OpenAI vision models &
\path{models/_openai_vision.py}; wrappers such as \path{models/gpt_5_5.py} &
\code{gpt-4.1}, \code{gpt-4o}, \code{gpt-5.*} API model IDs &
Hosted image-to-text transformer systems observed only through the Responses
API request and returned text. \\
Qwen2.5-VL &
\path{models/_qwen2_5_vl.py};
\path{models/qwen2_5_vl_3b_instruct.py} &
\path{Qwen/Qwen2.5-VL-3B-Instruct} &
Vision-language transformer using a Qwen processor, image inputs, chat
template, and autoregressive text decoding; optional PEFT adapter loading. \\
LLaVA &
\path{models/llava15_7b.py} &
\path{llava-hf/llava-1.5-7b-hf} &
CLIP-style visual tokens aligned with a causal language model through
image-token placeholders and checkpoint-specific processor settings. \\
PaliGemma &
\path{models/paligemma_3b_mix_224.py} &
\path{google/paligemma-3b-mix-224} &
PaLI/Gemma-style vision-language conditional generation with paired image
and text processing before decoder generation. \\
\bottomrule
\end{tabularx}
\normalsize
\end{table}

\begin{figure}[H]
\centering
\begin{tabularx}{0.95\textwidth}{>{\bfseries}p{29mm} X >{\centering\arraybackslash}p{32mm}}
\toprule
Path & Processing sequence visible in code & Output \\
\midrule
API vision wrapper &
Convert image to RGB JPEG data URL; pair \code{input\_text} with
\code{input\_image}; request a response from the configured model ID. &
Response output text \\
\addlinespace
Local generic wrapper &
Apply processor chat template to an image/text turn; build tensors; generate
tokens on the loaded model; decode only generated continuation. &
Decoded continuation \\
\addlinespace
Qwen2.5-VL &
Apply Qwen chat template and image processor; optionally attach a PEFT
adapter before generation. &
Decoded continuation \\
\addlinespace
LLaVA &
Prepare image-token prompt; synchronize image-token configuration when
necessary; generate and clean task-shaped text. &
Cleaned task answer \\
\bottomrule
\end{tabularx}
\caption{Operational paths exposed by the repository wrappers.}
\label{fig:paths}
\end{figure}

\section{The common multimodal transformer pattern}
The local open models in this report all use the same broad recipe.  An image
is split into patches or otherwise converted into visual tokens.  A vision
transformer turns those tokens into hidden states.  A small connector or
model-specific merger maps the visual states into the language model's token
space.  The language model is then used as an autoregressive decoder: it sees
the prompt plus visual context and generates answer tokens one at a time.

\begin{figure}[H]
\centering
\resizebox{\textwidth}{!}{%
\begin{tabular}{c c c c c c c}
\archbox{23mm}{Image\\pixels} & $\rightarrow$ &
\archbox{29mm}{Patchify /\\image processor} & $\rightarrow$ &
\archbox{31mm}{Vision\\Transformer} & $\rightarrow$ &
\archbox{32mm}{Visual token\\hidden states} \\
&&&&&& $\downarrow$ \\[-1mm]
\archbox{28mm}{Prompt\\tokens} & $\rightarrow$ &
\archbox{31mm}{Text\\embedding} & $\rightarrow$ &
\archbox{34mm}{Multimodal\\token stream} & $\rightarrow$ &
\archbox{34mm}{Decoder-only\\Transformer} \\
&&&&&& $\downarrow$ \\[-1mm]
&&&&&& \archbox{34mm}{Generated\\answer text}
\end{tabular}
}
\caption{Common architecture pattern for local vision-language transformers:
visual tokens are projected into a language-token stream before autoregressive
decoding.}
\label{fig:commonvlm}
\end{figure}

The differences between Qwen2.5-VL, LLaVA, and PaliGemma mostly lie in three
places: the visual encoder, the connector between vision and language, and the
language backbone.  Those choices affect how well a model handles detail,
layout, long prompts, and task-specific output formatting.

\section{API-served vision models}
The API pathway lives in \code{models/\_openai\_vision.py}.  It does very
little on purpose: convert the PIL image to RGB, encode it as a JPEG data URL,
send it with the cleaned prompt through the Responses API, and return the
output text.  The GPT-5 wrapper variant also avoids sending temperature where
that API configuration does not support it.

That gives the benchmark a convenient interface, but it also limits what we
can claim.  We can compare answer quality and observed request timing.  We
cannot inspect token-level visual features or provider-side memory use.  So
when the report compares API labels with local labels, it is comparing their
end-to-end behavior in this harness, not running a controlled hardware or
architecture ablation.

\begin{figure}[H]
\centering
\begin{tabular}{c c c c c}
\archbox{28mm}{PIL image} & $\rightarrow$ &
\archbox{35mm}{RGB JPEG\\data URL} & $\rightarrow$ &
\archbox{42mm}{Responses API\\\code{input\_image}} \\
&&&& $\downarrow$ \\[-1mm]
\archbox{28mm}{Prompt text} & $\rightarrow$ &
\archbox{35mm}{Cleaned\\instruction} & $\rightarrow$ &
\archbox{42mm}{Responses API\\\code{input\_text}} \\
&&&& $\downarrow$ \\[-1mm]
&&&& \archbox{42mm}{Returned\\output text}
\end{tabular}
\caption{Observable architecture boundary for API-served model labels.  The
repo can describe the request and response path, but not the provider-side
transformer internals.}
\label{fig:openaiarch}
\end{figure}

\section{Local Hugging Face models}
The local wrappers share loading code from \code{models/\_hf\_model.py}.  They
try the local cache first and fall back to the hub when needed.  The generic
image-text route builds a one-image, one-instruction conversation, applies the
processor's chat template, moves tensors to the model device, generates a
continuation, and decodes only the new tokens.

This route is useful because we can read and adjust the exact prompt
preparation, generation settings, and local telemetry.  The trade-off is
practical: local hardware decides which checkpoints are feasible and how slow
the run will be.

\begin{figure}[H]
\centering
\resizebox{\textwidth}{!}{%
\begin{tabular}{c c c c c c c}
\archbox{25mm}{Model ID} & $\rightarrow$ &
\archbox{32mm}{Processor\\load} & $\rightarrow$ &
\archbox{32mm}{Model\\load} & $\rightarrow$ &
\archbox{32mm}{Device map\\and dtype} \\
&&&&&& $\downarrow$ \\[-1mm]
\archbox{25mm}{Image} & $\rightarrow$ &
\archbox{32mm}{Chat\\template} & $\rightarrow$ &
\archbox{32mm}{Tensor\\batch} & $\rightarrow$ &
\archbox{32mm}{Generate\\tokens}
\end{tabular}
}
\caption{Generic Hugging Face execution graph used by the local wrappers:
processor and model loading are separate from the per-example generation path.}
\label{fig:hfarch}
\end{figure}

\section{Qwen2.5-VL and adapter loading}
The Qwen2.5-VL wrapper uses
\code{Qwen2\_5\_VLForConditionalGeneration}.  It builds a user conversation
with an image placeholder and the benchmark prompt, lets the processor prepare
the text and image inputs, generates a bounded continuation, and decodes the
answer.  If an adapter path is provided, the wrapper attaches a PEFT adapter
before evaluation.

That adapter hook matters because it lets the report move from comparison to
intervention.  The baseline runs ask how the model behaves as-is.  The
Fashion-MNIST case study asks whether a small task-specific update changes
that behavior, while keeping the benchmark runner and answer parsing the same.

Architecturally, Qwen2.5-VL follows a three-part vision-language transformer
design.  The image side uses a native dynamic-resolution Vision Transformer,
so the processor can preserve more of the image's original spatial structure
instead of forcing every input into one fixed token pattern.  Window attention
keeps the vision encoder more efficient on high-resolution inputs.  A
multimodal merger then maps visual representations into the Qwen language
model's embedding space.  The decoder is a Qwen transformer language backbone,
with multimodal rotary position handling used to keep visual and textual
positions aligned.

\begin{figure}[H]
\centering
\begin{tabular}{c c c c c}
\archbox{31mm}{Native-resolution\\image} & $\rightarrow$ &
\archbox{36mm}{Qwen image\\processor} & $\rightarrow$ &
\archbox{38mm}{Dynamic-resolution\\ViT} \\
&&&& $\downarrow$ \\[-1mm]
\archbox{31mm}{Prompt +\\chat template} & $\rightarrow$ &
\archbox{36mm}{Text tokens} & $\rightarrow$ &
\archbox{38mm}{Multimodal\\merger / MRoPE} \\
&&&& $\downarrow$ \\[-1mm]
&&&& \archbox{38mm}{Qwen decoder\\Transformer} \\
&&&& $\downarrow$ \\[-1mm]
&&&& \archbox{38mm}{Answer tokens}
\end{tabular}
\caption{Qwen2.5-VL architecture as used by the local wrapper: a dynamic-
resolution ViT feeds a multimodal token stream consumed by the Qwen decoder.}
\label{fig:qwenarch}
\end{figure}

The adapter path adds a small trainable layer on top of that base architecture.
In the fine-tuning experiment, LoRA modules are attached to attention
projection matrices.  During evaluation the benchmark still sees the same
\code{predict(image, prompt)} interface; the difference is that the Qwen
decoder now contains the loaded PEFT adapter.

\begin{figure}[H]
\centering
\begin{tabular}{c c c c c}
\archbox{34mm}{Base Qwen2.5-VL\\weights} & $+$ &
\archbox{34mm}{LoRA adapter\\weights} & $\rightarrow$ &
\archbox{40mm}{Adapted generation\\path} \\
&&&& $\downarrow$ \\[-1mm]
\archbox{34mm}{Same benchmark\\prompt} & $\rightarrow$ &
\archbox{34mm}{Same parser\\and labels} & $\rightarrow$ &
\archbox{40mm}{Comparable paired\\predictions}
\end{tabular}
\caption{Adapter loading graph for the Qwen fine-tuning case study.  The model
changes, but the benchmark input and scoring path stay fixed.}
\label{fig:qwenadapterarch}
\end{figure}

\section{LLaVA's practical compatibility handling}
The LLaVA wrapper is a reminder that a checkpoint ID is not enough.  Some
LLaVA-style checkpoints need the processor settings to match the visual token
sequence expected by the model.  The wrapper reads the visual patch settings
and repairs image-token mismatches when needed.  It also cleans the generated
text differently by task: a classification answer is reduced to one line,
while detection output has to keep parseable labelled boxes.

This is not just housekeeping.  A bad prompt template or image-token mismatch
can make a capable checkpoint look broken.  The wrapper is therefore part of
the experiment, not just plumbing around it.

LLaVA-1.5 is a deliberately simple multimodal transformer stack.  A CLIP
Vision Transformer encodes the image into patch-level visual states.  A
vision-language connector, implemented as an MLP in LLaVA-1.5, projects those
visual states into the language model's embedding space.  The prompt contains
an image placeholder, and the projected visual tokens occupy that position in
the sequence seen by the language model.  The answer is generated by an
autoregressive decoder language model.

\begin{figure}[H]
\centering
\begin{tabular}{c c c c c}
\archbox{30mm}{Image} & $\rightarrow$ &
\archbox{35mm}{CLIP ViT\\vision encoder} & $\rightarrow$ &
\archbox{35mm}{Patch-level\\visual states} \\
&&&& $\downarrow$ \\[-1mm]
\archbox{30mm}{Prompt with\\\code{<image>}} & $\rightarrow$ &
\archbox{35mm}{Tokenizer and\\text embeddings} & $\rightarrow$ &
\archbox{35mm}{MLP projector\\+ image tokens} \\
&&&& $\downarrow$ \\[-1mm]
&&&& \archbox{35mm}{Causal LM\\decoder} \\
&&&& $\downarrow$ \\[-1mm]
&&&& \archbox{35mm}{Cleaned task\\answer}
\end{tabular}
\caption{LLaVA-style architecture: a CLIP vision transformer is connected to a
causal language model through a learned projector.}
\label{fig:llavaarch}
\end{figure}

The evaluated \code{llava-gemma-2b} label follows the same broad idea---vision
tokens connected to a decoder language model---but with a different language
backbone and checkpoint configuration.  That is why the report treats it as a
LLaVA-style local model while still relying on the wrapper's compatibility
logic for the exact processor settings.

\section{PaliGemma and compatible processing}
The PaliGemma wrapper chooses the available compatible conditional-generation
class, processes the image and prompt together, and decodes the generated
tokens.  The path is simpler than the LLaVA repair logic, but the results are
one of the more useful reminders in the report: a model can be weak in the
aggregate table and still be competitive for a particular family, in this case
QA.

PaliGemma is built from open components in the PaLI/Gemma line: a SigLIP vision
encoder and a Gemma language model joined for conditional generation.  Compared
with the LLaVA wrapper, the local code path is less about repairing an explicit
\code{<image>} placeholder and more about giving the processor the image and
text together so the model can form the combined sequence expected by the
checkpoint.

\begin{figure}[H]
\centering
\begin{tabular}{c c c c c}
\archbox{30mm}{Image} & $\rightarrow$ &
\archbox{35mm}{SigLIP-style\\vision encoder} & $\rightarrow$ &
\archbox{35mm}{Visual\\representations} \\
&&&& $\downarrow$ \\[-1mm]
\archbox{30mm}{Prompt} & $\rightarrow$ &
\archbox{35mm}{PaliGemma\\processor} & $\rightarrow$ &
\archbox{35mm}{Joint image-text\\inputs} \\
&&&& $\downarrow$ \\[-1mm]
&&&& \archbox{35mm}{Gemma decoder\\generation} \\
&&&& $\downarrow$ \\[-1mm]
&&&& \archbox{35mm}{Decoded\\answer text}
\end{tabular}
\caption{PaliGemma architecture: SigLIP-style visual representations are paired
with a Gemma language decoder for image-conditioned generation.}
\label{fig:paligemmaarch}
\end{figure}

\section{Where the architectures differ}
The practical differences between these models can be summarized by where
vision enters the transformer computation.  In the API-served path, the
repository sees only the boundary.  In Qwen2.5-VL, the model-specific processor
and merger create a native multimodal sequence for the Qwen decoder.  In
LLaVA, the connector explicitly projects CLIP visual states into the language
embedding space.  In PaliGemma, the processor builds the paired image-text
inputs expected by the SigLIP-plus-Gemma conditional-generation stack.

\begin{figure}[H]
\centering
\begin{tabularx}{0.98\textwidth}{>{\bfseries}p{24mm} X X X}
\toprule
Family & Vision side & Fusion point & Language side \\
\midrule
OpenAI API path &
Not inspectable in this repository. &
Provider-side, not visible to the harness. &
Returned text only. \\
Qwen2.5-VL &
Dynamic-resolution ViT with efficient visual attention. &
Model-specific multimodal merger and position handling. &
Qwen decoder transformer. \\
LLaVA &
CLIP Vision Transformer visual states. &
MLP projector inserts visual tokens into the language stream. &
Causal language model decoder. \\
PaliGemma &
SigLIP-style vision encoder. &
Processor/model build paired image-text inputs. &
Gemma decoder transformer. \\
\bottomrule
\end{tabularx}
\caption{Comparison graph for the main architecture families.  The central
difference is the fusion point between vision features and language decoding.}
\label{fig:archcompare}
\end{figure}

\resultbox{\textbf{Operational conclusion.} The evaluation treats each system
as an image-to-text component.  The wrappers handle the family-specific details
needed to avoid obvious formatting failures.  The results are therefore best
read as performance of a model \emph{plus its wrapper and prompt protocol}.}

\clearpage
\chapter{Datasets and Benchmarks}

\section{Twenty-one benchmarks, four families}
The aggregate run contains twenty-one benchmarks, grouped into captioning, QA,
labeling, and detection.  These families ask for different kinds of text, so
they reward different habits in a model.  Captioning needs a short description.
QA needs a concise answer to a visual question.  Labeling asks the model to
choose from a class list.  Detection asks for both a label and usable
coordinates.

\begin{table}[H]
\centering
\caption{Task families represented in the current aggregate comparison.}
\begin{tabularx}{\textwidth}{>{\bfseries}p{27mm} p{55mm} X}
\toprule
Family & Examples appearing in output & What the model must return \\
\midrule
Captioning & \code{flickr30k}, \code{textcaps},
\code{mscoco\_caption}, \code{conceptual\_captions\_caption} &
A short natural-language description. \\
QA & \code{docvqa}, \code{gqa}, \code{vqav2} &
A question-conditioned answer string. \\
Labeling & \code{fashion\_mnist}, \code{imagenet1k}, \code{places},
\code{inaturalist}, \code{lsun}, \code{mvtec\_ad}, video/image label tasks &
One label from a candidate list. \\
Detection & \code{mscoco}, \code{openimages\_v4\_detection},
\code{flickr30k\_entities} &
Labels paired with parseable boxes. \\
\bottomrule
\end{tabularx}
\end{table}

\section{Coverage is broad enough to compare, not to generalize universally}
Each of the fourteen model labels appears on each of the twenty-one current
benchmarks, giving a rectangular 294-cell comparison.  That is important:
the aggregate table is not inflated by one model being evaluated only on its
best tasks.  Still, this is not a universal ranking of multimodal systems.
Some task runs have modest sample sizes, and several benchmarks produce large
ties.

\begin{figure}[H]
\centering
\includegraphics[width=0.9\textwidth]{\figpath{telemetry_coverage.png}}
\caption{Coverage of recorded telemetry in the main comparison.  The
regenerated output marks all evaluated cells as measured rather than
estimated.}
\label{fig:coverage}
\end{figure}

\section{Reading the figures together}
The results section uses three views repeatedly.  The normalized score plot is
the shortest answer to ``who did well across these mixed tasks?''  The raw
score plot keeps the original metric scale visible, but should not carry the
whole comparison by itself.  Heatmaps and family plots show where the averages
come from.  When these views disagree, that disagreement is usually the point,
not a problem to smooth away.

\clearpage
\chapter{Evaluation Methodology}

\section{The evaluation loop}
The central pipeline is deliberately plain.  A dataset provides a row and an
image.  A benchmark turns that row into a constrained prompt.  A model wrapper
runs its own inference path and returns text.  The benchmark parses that text,
records whether it was useful, and stores telemetry.  The comparison script
then reads the JSON reports and builds the tables and plots used here.

\begin{figure}[H]
\centering
\setlength{\tabcolsep}{5pt}
\begin{tabular}{>{\centering\arraybackslash}p{25mm} c
                >{\centering\arraybackslash}p{30mm} c
                >{\centering\arraybackslash}p{29mm} c
                >{\centering\arraybackslash}p{29mm}}
\fcolorbox{accent}{soft}{\parbox{22mm}{\centering Dataset\\row + image}} &
$\rightarrow$ &
\fcolorbox{accent}{soft}{\parbox{27mm}{\centering Benchmark\\prompt}} &
$\rightarrow$ &
\fcolorbox{accent}{soft}{\parbox{26mm}{\centering Model\\text output}} &
$\rightarrow$ &
\fcolorbox{accent}{soft}{\parbox{26mm}{\centering Score\\+ telemetry}}\\
\end{tabular}
\caption{The observable evaluation pipeline shared by the model wrappers.}
\label{fig:pipeline}
\end{figure}

For labeling tasks, the prompt asks for exactly one label from a supplied
list.  The benchmark normalizes the model's text and counts it as correct if
it matches one of the valid labels for that row.  When the label set is large,
the benchmark keeps the valid labels and samples distractors with a row-derived
seed.  That keeps the choice set stable rather than adding avoidable noise.

Captioning, QA, and detection keep the same input/output shape but use
different scoring logic.  Captioning scores caption text.  QA scores answer
text.  Detection parses labelled coordinate expressions.  Some video-oriented
datasets are represented as contact sheets of frames, which keeps the model
contract image-based while still giving the model temporal clues.

\section{What is recorded}
Each JSON report contains per-example outputs and aggregate statistics.  The
comparison builder extracts the primary score, result count, wall-clock time,
generation time, memory peaks, and whether telemetry was measured.  In the
regenerated table used here, all 294 main cells have measured telemetry; none
are marked as estimated.

\begin{table}[H]
\centering
\caption{Scoring concepts used in the comparison builder.}
\begin{tabularx}{\textwidth}{p{38mm} X}
\toprule
Quantity & Meaning in this report \\
\midrule
Raw primary score & Original score for a benchmark, useful within its task
type but not automatically commensurate across families. \\
Within-benchmark normalized score & For each benchmark, a model's position
between the observed minimum and maximum; used for the overall ranking. \\
Benchmark rank & Rank of each model on a benchmark, with tied maxima receiving
the same first rank. \\
Pairwise win rate & Fraction of shared benchmarks on which model A beats
model B, with ties contributing half in the generated summary. \\
Bootstrap interval & Variation in mean raw score under benchmark resampling;
used as a caution about ordering stability. \\
\bottomrule
\end{tabularx}
\end{table}

\section{A necessary warning about normalization}
Within-benchmark normalization asks a useful question: compared with the other
models on this same task, how strong was this system?  That prevents metrics
with naturally larger numbers from dominating the aggregate.  The trade-off is
that a small absolute difference on a tightly clustered benchmark can look more
important than it really is, while an all-tie benchmark adds little evidence.
For that reason, the report keeps normalized rankings next to raw plots,
spreads, ranks, and family-level discussion.

\clearpage
\chapter{Baseline Results and Comparison}

\section{The leading aggregate result}
The normalized comparison has a clear leader: \code{gpt-5.5} reaches 0.825
with a mean benchmark rank of 2.76.  \code{gpt-4.1} follows at 0.757, a gap
of 0.068 normalized points.  Behind them, \code{gpt-5.1},
\code{gpt-5.4-mini}, \code{gpt-5.4}, and
\code{qwen2.5-vl-3b-instruct} sit in a fairly tight group from 0.742 down to
0.714.  I would not read every adjacent rank in that group as a decisive
practical difference.

\begin{figure}[H]
\centering
\includegraphics[width=0.92\textwidth]{\figpath{model_normalized_scores.png}}
\caption{Mean score normalized within each benchmark.  This is the primary
cross-family ranking used in the report.}
\label{fig:normalized}
\end{figure}

\begin{table}[H]
\centering
\caption{Top ten models in the aggregate output.}
\begin{tabular}{lrrrr}
\toprule
Model & Normalized mean & Raw mean & Mean rank & Firsts \\
\midrule
gpt-5.5 & 0.825 & 0.555 & 2.76 & 8 \\
gpt-4.1 & 0.757 & 0.518 & 3.48 & 9 \\
gpt-5.1 & 0.742 & 0.508 & 3.76 & 9 \\
gpt-5.4-mini & 0.727 & 0.496 & 4.48 & 7 \\
gpt-5.4 & 0.724 & 0.505 & 4.81 & 6 \\
qwen2.5-vl-3b-instruct & 0.714 & 0.478 & 4.81 & 10 \\
gpt-4o & 0.708 & 0.495 & 4.62 & 6 \\
gpt-5 & 0.688 & 0.499 & 4.81 & 8 \\
gpt-5.2 & 0.678 & 0.480 & 5.19 & 6 \\
llava-1.5-7b-hf & 0.666 & 0.467 & 4.57 & 7 \\
\bottomrule
\end{tabular}
\end{table}

\section{Raw scores give context, not a replacement leaderboard}
Raw means tell a similar top-line story: \code{gpt-5.5} leads at 0.555,
compared with 0.518 for \code{gpt-4.1}.  That agreement is reassuring, but it
does not make raw averaging enough.  A classification task with values near
one and a detection task with values near zero do not carry equal weight in a
raw average.  The raw chart is best used as a reality check before moving into
the family-level analysis.

\section{Winning often is not the same as winning broadly}
The first-place count shows the most interesting tension in the table.
Qwen2.5-VL 3B wins or ties for first on ten benchmarks, more than any other
label, but it ranks sixth by normalized mean.  In other words, it has real
peaks and real weak spots.  \code{gpt-5.5} has fewer first-place finishes,
but is steadier across the full set.

\begin{figure}[H]
\centering
\begin{minipage}{0.49\textwidth}
\centering
\includegraphics[width=\linewidth]{\figpath{model_mean_scores.png}}\\
{\small (a) Mean raw primary score.}
\end{minipage}\hfill
\begin{minipage}{0.49\textwidth}
\centering
\includegraphics[width=\linewidth]{\figpath{first_place_finishes.png}}\\
{\small (b) First-place finishes including ties.}
\end{minipage}
\caption{Two complementary aggregate views.  Raw scores preserve original
metric scale, while win counts reveal Qwen's task-specific peaks.}
\label{fig:rawfirsts}
\end{figure}

\clearpage
\section{Task-family comparison}

\subsection{Four families tell four different stories}
The overall score is useful for a quick ranking, but the models make more
sense once the tasks are separated.  The family summary produces a different
leader in almost every category.  That is not noise; it tells us that captions,
short answers, labels, and boxes stress the models in different ways.

\begin{table}[H]
\centering
\caption{Family leaders by benchmark-normalized mean score.}
\begin{tabularx}{\textwidth}{>{\bfseries}p{26mm} X r r}
\toprule
Family & Leader or tied leaders & Normalized & Raw mean \\
\midrule
Captioning & qwen2.5-vl-3b-instruct & 0.977 & 0.325 \\
Detection & gpt-5.4; gpt-5.5 & 0.873 & 0.159 \\
Labeling & gpt-4.1; gpt-5.1 & 0.877 & 0.764 \\
QA & paligemma-3b-mix-224 & 0.833 & 0.600 \\
\bottomrule
\end{tabularx}
\end{table}

\begin{figure}[H]
\centering
\includegraphics[width=0.98\textwidth]{\figpath{family_small_multiples.png}}
\caption{Scores separated into task families and component benchmarks.  The
panel view prevents a family average from concealing its source.}
\label{fig:families}
\end{figure}

\subsection{Captioning: local Qwen and LLaVA are competitive for a reason}
Captioning is where Qwen looks strongest.  \code{qwen2.5-vl-3b-instruct}
gets a normalized family mean of 0.977.  It wins three captioning rows:
\code{conceptual\_captions\_caption}, \code{flickr30k}, and \code{textcaps}.
\code{llava-1.5-7b-hf} is close behind at 0.936 and wins
\code{mscoco\_caption}.  Their raw captioning means are almost identical after
rounding, but normalization favors Qwen because it stays closer to the best
result within the individual caption tasks.

For an application mostly built around descriptions, this matters more than
the aggregate ranking.  Picking Qwen for captioning would not be a preference
for a lower-ranked model; it would be following the more relevant evidence.

\subsection{QA and labeling: an instructive reversal}
PaliGemma is the best warning against dismissing a model from one aggregate
number.  It has the lowest normalized aggregate mean in the main table, 0.318,
but it leads the normalized QA family summary at 0.833.  At the same time,
\code{gpt-5.5} has the higher raw QA mean, 0.633 versus PaliGemma's 0.600,
while landing behind it on normalized QA mean at 0.800.

That reversal is not a contradiction.  The normalized summary gives each QA
benchmark a relative voice; the raw average lets the original numeric scales
matter more.  Both facts belong in the report: \code{gpt-5.5} has the higher
raw QA average here, while PaliGemma is the stronger relative competitor
across the individual QA tasks after normalization.

Labeling is less surprising, but the top is tight.  \code{gpt-4.1} and
\code{gpt-5.1} tie for the normalized family lead at 0.877 and the raw family
lead at 0.764.  \code{gpt-5.5} follows closely at 0.864 normalized.  For a
label-heavy deployment, latency and operating constraints may matter more than
these small score differences.

\begin{figure}[H]
\centering
\includegraphics[width=0.75\textwidth]{\figpath{family_radar.png}}
\caption{Family profile view.  The chart is useful for seeing specialization,
but exact conclusions should be read from the accompanying family table.}
\label{fig:radar}
\end{figure}

\subsection{Detection: low values and a non-discriminating task}
Detection sits on the lowest raw score scale of the four families.  In the
family mean, \code{gpt-5.4} and \code{gpt-5.5} tie at 0.873 normalized
(raw 0.159), with \code{gpt-5.4-mini} close behind at 0.806 (raw 0.154).
But this section needs a warning: \code{flickr30k\_entities} gives every
evaluated model a raw score of zero.  That does not prove the models are
equally bad at detection.  It says this benchmark/protocol combination did not
distinguish them on that row.

The other detection rows are more informative.  On \code{mscoco},
\code{gpt-4.1} and \code{gpt-5.1} tie at 0.114.  On
\code{openimages\_v4\_detection}, \code{gpt-5.4-mini} leads at 0.413, just
ahead of \code{gpt-5.4} and \code{gpt-5.5}, both at 0.405.  Detection belongs
in the report, but the all-zero row should be debugged before this part of the
comparison is treated as settled.

\begin{figure}[H]
\centering
\includegraphics[width=0.99\textwidth]{\figpath{score_heatmap_normalized.png}}
\caption{Normalized score heatmap.  Bright cells expose specialist wins;
flat or uniformly dark rows expose weak discriminatory evidence.}
\label{fig:heatmap}
\end{figure}

\clearpage
\section{Spread, ranks, and uncertainty}

\subsection{Which benchmarks separate models?}
Three benchmarks---\code{inaturalist}, \code{lsun}, and \code{ucf101}---span
the complete observed range from 0.000 to 1.000.  \code{docvqa} follows with
a spread of 0.800, while \code{fashion\_mnist} and \code{mvtec\_ad} each
span 0.700.  These rows pull the eye in a heatmap, so they deserve extra
attention.  Before treating them as broad capability signals, it is worth
checking the example counts and the actual mistakes.

\subsection{Ranks retain ties that averages obscure}
Rank views help where the score scales differ.  They also expose tie-heavy
rows, including tasks where many models reach the maximum.  A first-place
count should therefore be read as ``reached the best recorded value,'' not
``uniquely solved the task.''

\begin{figure}[H]
\centering
\begin{minipage}{0.47\textwidth}
\centering
\includegraphics[width=\linewidth]{\figpath{benchmark_spread.png}}\\
{\small (a) Observed benchmark spread.}
\end{minipage}\hfill
\begin{minipage}{0.51\textwidth}
\centering
\includegraphics[width=\linewidth]{\figpath{rank_heatmap.png}}\\
{\small (b) Rank heatmap, including ties.}
\end{minipage}
\caption{Spread and rank should be read together: separating tasks can also
contain substantial ties or limited evidence.}
\label{fig:spreadranks}
\end{figure}

\subsection{Head-to-head comparisons are close at the top}
The aggregate leader does not sweep its closest competitors.  \code{gpt-5.5}
beats \code{gpt-4.1} on seven benchmarks, loses on five, and ties on nine;
its mean raw difference is only +0.037.  Against
\code{qwen2.5-vl-3b-instruct}, it wins eight, loses six, and ties seven, with
a mean difference of +0.077.  Those are real aggregate advantages, but not
evidence that either alternative is generally unsuitable.

The bigger pairwise gaps appear farther down the table.  \code{gpt-5.5}
records 17 wins, no losses, and four ties against \code{gpt-5.4-nano}; against
\code{paligemma-3b-mix-224} it records 18 wins, two losses, and one tie.  Even
that last comparison should be read alongside PaliGemma's QA result.

\subsection{Intervals make overstatement harder}
The bootstrap intervals are generated by resampling benchmarks, not individual
examples.  They show how sensitive the raw mean is to the particular task mix
in this report.  \code{gpt-5.5} has observed raw mean 0.555 with a 90 percent
interval of [0.442, 0.665].  \code{gpt-4.1} is 0.518 with interval
[0.401, 0.639], and \code{gpt-5.1} is 0.508 with [0.394, 0.631].  Those
intervals overlap a lot.

\begin{figure}[H]
\centering
\begin{minipage}{0.49\textwidth}
\centering
\includegraphics[width=\linewidth]{\figpath{pairwise_win_rate_heatmap.png}}\\
{\small (a) Pairwise win-rate heatmap.}
\end{minipage}\hfill
\begin{minipage}{0.49\textwidth}
\centering
\includegraphics[width=\linewidth]{\figpath{bootstrap_confidence_intervals.png}}\\
{\small (b) Bootstrap intervals for raw means.}
\end{minipage}
\caption{The mixed head-to-head pattern and overlapping intervals both
caution against treating the top ordering as universal.}
\label{fig:pairwisebootstrap}
\end{figure}

The fair reading is balanced.  \code{gpt-5.5} leads this regenerated table by
both raw and normalized summaries.  At the same time, the result set is not
large or uniform enough to promise the same order under a different application
mix.  In practice, the family plots and latency requirements should often make
the final decision.

\section{Runtime and resource trade-offs}

\subsection{Recorded timing}
All main comparison cells report measured timing rather than estimated
statistics.  The fastest mean wall-clock values are \code{gpt-5.4-nano}
(1.18 seconds), \code{llava-1.5-7b-hf} (1.19 seconds), and
\code{gpt-5.4-mini} (1.21 seconds).  The aggregate quality leader,
\code{gpt-5.5}, averages 1.63 seconds.  \code{qwen2.5-vl-3b-instruct}
averages 2.70 seconds.  PaliGemma and LLaVA-Gemma are much slower in this
recorded table, at 12.81 and 99.07 seconds respectively.

\begin{figure}[H]
\centering
\begin{minipage}{0.49\textwidth}
\centering
\includegraphics[width=\linewidth]{\figpath{telemetry_mean_wall_clock_seconds.png}}\\
{\small (a) Mean wall-clock seconds.}
\end{minipage}\hfill
\begin{minipage}{0.49\textwidth}
\centering
\includegraphics[width=\linewidth]{\figpath{telemetry_peak_cpu_ram_gib.png}}\\
{\small (b) Recorded peak CPU RAM.}
\end{minipage}
\caption{Runtime and memory measurements observed by the benchmark harness.
Absolute values combine model pathway and execution environment.}
\label{fig:wallclockram}
\end{figure}

\subsection{A practical quality--latency frontier}
If we prefer higher normalized score and lower wall-clock time, the observed
non-dominated frontier contains five labels: \code{gpt-5.4-nano},
\code{llava-1.5-7b-hf}, \code{gpt-5.4-mini}, \code{gpt-4.1}, and
\code{gpt-5.5}.  Moving along that frontier buys more aggregate quality for
modest increases in recorded time.

Qwen is not on that aggregate frontier because its 0.714 normalized score and
2.70-second mean are dominated in the all-task summary.  It can still be the
right choice for a caption-heavy application.  The frontier uses the aggregate
objective; Qwen's strongest evidence is family-specific.

\subsection{Memory values require environmental context}
The memory chart is most useful for local runs, where model size and device
constraints decide what is feasible.  It is less useful for saying that one
API-served system is intrinsically lighter than another.  The repository
measures its own process and observed call, not the provider's serving
infrastructure.  The same caution applies to direct API-versus-local timing.

\begin{figure}[H]
\centering
\includegraphics[width=0.98\textwidth]{\figpath{telemetry_wall_clock_by_benchmark.png}}
\caption{Wall-clock results separated by benchmark.  Per-task variation can
matter as much as a model mean when designing an interactive workflow.}
\label{fig:wallbybench}
\end{figure}

\resultbox{\textbf{Deployment implication.} For a mixed-task workload in this
measurement context, \code{gpt-5.5} gives the strongest score with moderate
recorded time.  For captioning, Qwen remains compelling despite slower
aggregate timing.  A production latency decision still needs repeated runs in
the intended local or API environment.}

\clearpage
\chapter{Fine-Tuning Experiments}

\section{Why a separate case study?}
The main comparison asks how ready-to-use wrappers behave across task
families.  Fine-tuning asks a different question: can we change one model so it
handles a narrow task better?  For that reason, the fine-tuning result files
stay outside the aggregate leaderboard and are analyzed as paired Fashion-MNIST
evaluations.

Fashion-MNIST is a useful pilot because it is small, visual, and has ten clear
classes.  It also lets training data come from the training split while
evaluation stays on the held-out test split.  At the same time, it is not as
easy as it looks for a language-generating vision model.  Classes such as
\emph{shirt}, \emph{pullover}, \emph{coat}, and \emph{t-shirt/top} are close
enough visually that class collapse appears quickly.

\section{QLoRA workflow}
The generic script in \code{fine\_tuning/train\_qwen25vl\_benchmark.py} reads
manifest records with images, prompts, and target answers.  The collator
formats each example as a Qwen conversation and masks the prompt tokens so the
loss is applied to the assistant answer rather than the instruction.  The base
model is loaded with 4-bit NF4 quantization and double quantization.  LoRA
modules are attached to the attention projections \code{q\_proj},
\code{k\_proj}, \code{v\_proj}, and \code{o\_proj}.  The standard Qwen wrapper
then loads the adapter for evaluation.

\begin{figure}[H]
\centering
\begin{tabular}{c c c c c}
\fcolorbox{accent}{soft}{\parbox{24mm}{\centering Train split\\images}} &
$\rightarrow$ &
\fcolorbox{accent}{soft}{\parbox{30mm}{\centering Balanced\\manifests}} &
$\rightarrow$ &
\fcolorbox{accent}{soft}{\parbox{33mm}{\centering QLoRA\\adapter}} \\
&&&& $\downarrow$ \\[-1mm]
\fcolorbox{accent}{soft}{\parbox{29mm}{\centering Test split\\only}} &
$\rightarrow$ &
\fcolorbox{accent}{soft}{\parbox{31mm}{\centering Same benchmark\\runner}} &
$\rightarrow$ &
\fcolorbox{accent}{soft}{\parbox{33mm}{\centering Paired\\errors}} \\
\end{tabular}
\caption{Fine-tuning and held-out evaluation workflow.  The test examples do
not enter the training manifests.}
\label{fig:ftpipe}
\end{figure}

\section{Accuracy improves, but the error distribution moves}
The earlier LoRA run matches the base result exactly: 19 correct answers out
of 50.  The balanced LoRA run reaches 24 correct answers, a gain of five
answers, or ten percentage points, on the same paired sample.

\begin{table}[H]
\centering
\caption{Stored Fashion-MNIST Qwen2.5-VL-3B evaluation results.}
\begin{tabular}{lrrr}
\toprule
Run & Correct & Total & Accuracy \\
\midrule
Base test-split evaluation & 19 & 50 & 38.0\% \\
Earlier LoRA adapter & 19 & 50 & 38.0\% \\
Balanced LoRA adapter & 24 & 50 & 48.0\% \\
\bottomrule
\end{tabular}
\end{table}

\begin{figure}[H]
\centering
\begin{tabular}{@{}l@{\hspace{8mm}}l r@{}}
Base & \textcolor{accent}{\rule{57mm}{7pt}} & 38.0\% \\
Earlier LoRA & \textcolor{accent}{\rule{57mm}{7pt}} & 38.0\% \\
Balanced LoRA & \textcolor{accent}{\rule{72mm}{7pt}} & 48.0\% \\
\end{tabular}
\caption{Accuracy comparison for the paired 50-example Fashion-MNIST
evaluation.}
\label{fig:ftbars}
\end{figure}

The paired transition count says more than the bar chart.  Compared with the
base predictions, the balanced adapter changes 17 answers.  Ten of those
changes fix a base error.  Five turn a previously correct answer into an
error.  The remaining changed answers move from one wrong label to another.
So the adapter helps, but it does not improve the model uniformly.

\begin{table}[H]
\centering
\caption{Correctness transitions from base to balanced adapter.}
\begin{tabular}{lr}
\toprule
Transition & Count \\
\midrule
Incorrect under base, correct under balanced adapter & 10 \\
Correct under base, incorrect under balanced adapter & 5 \\
Net gain in correct predictions & +5 \\
All changed predictions & 17 \\
\bottomrule
\end{tabular}
\end{table}

\section{Class-level improvement and class bias}
The adapter improves \emph{shirt} sharply, from 2/5 correct examples to 5/5.
It also fixes some \emph{coat}, \emph{trouser}, \emph{ankle boot}, \emph{bag},
and \emph{sandal} examples.  But the gain has a cost.  \emph{Pullover} drops
from 4/6 to 1/6, and \emph{t-shirt/top} drops from 2/3 to 0/3.  Those losses
are not random-looking: they sit in the same visually similar clothing region
as the newly strengthened shirt label.

\begin{table}[H]
\centering
\caption{Correct predictions by class: base versus balanced LoRA.}
\begin{tabular}{lrr@{\hspace{7mm}}lrr}
\toprule
Class & Base & Balanced & Class & Base & Balanced \\
\midrule
ankle boot & 1/4 & 2/4 & sandal & 0/4 & 1/4 \\
bag & 2/4 & 3/4 & shirt & 2/5 & 5/5 \\
coat & 0/5 & 2/5 & sneaker & 7/7 & 7/7 \\
dress & 1/5 & 1/5 & t-shirt/top & 2/3 & 0/3 \\
pullover & 4/6 & 1/6 & trouser & 0/7 & 2/7 \\
\bottomrule
\end{tabular}
\end{table}

\begin{figure}[H]
\centering
\begin{tabular}{@{}l@{\hspace{3mm}}l@{\hspace{5mm}}l@{\hspace{3mm}}l@{}}
\multicolumn{2}{c}{\textbf{Base predictions}} &
\multicolumn{2}{c}{\textbf{Balanced predictions}}\\[4pt]
shirt & \textcolor{warning}{\rule{50mm}{6pt}} 20 &
shirt & \textcolor{warning}{\rule{59mm}{6pt}} 25 \\
sneaker & \textcolor{accent}{\rule{25mm}{6pt}} 10 &
sneaker & \textcolor{accent}{\rule{20mm}{6pt}} 8 \\
pullover & \textcolor{accent}{\rule{22.5mm}{6pt}} 9 &
coat & \textcolor{accent}{\rule{12.5mm}{6pt}} 5 \\
\end{tabular}
\caption{Most frequent predicted labels.  The balanced adapter improves
accuracy while increasing an already prominent preference for \emph{shirt}.}
\label{fig:bias}
\end{figure}

There are only five true shirt instances in the fifty-example evaluation, yet
the balanced adapter predicts shirt twenty-five times.  That is the main
caution in the case study.  Balanced training improves measured accuracy, but
it does not solve calibration among related garment classes.  Before treating
this adapter as more than a pilot, we would need a larger test, a confusion
matrix, and repeated training seeds.

\section{The Colab smoke result}
A newer downloaded artifact contains a twenty-example base-versus-adapter
smoke evaluation from the generalized Colab workflow.  The base run gets 7/20
(35.0 percent), and the adapter run gets 9/20 (45.0 percent).  Ten predictions
change, five errors are fixed, and three correct base outputs regress.

That is useful operational evidence: the generalized training and evaluation
path can produce an adapter, load it, and score it with the same runner.  It is
not strong additional performance evidence.  Twenty examples are too few for a
stable accuracy estimate, and this result should not be merged with the
fifty-example comparison as if it were a planned replication.

\begin{table}[H]
\centering
\caption{Interpretation of the two fine-tuning evidence sources.}
\begin{tabularx}{\textwidth}{p{35mm} p{31mm} X}
\toprule
Evidence source & Result & Suitable conclusion \\
\midrule
Paired 50-example stored comparison &
38.0\% to 48.0\% &
Balanced LoRA improved this held-out pilot sample, with visible class bias. \\
Downloaded 20-example Colab smoke pair &
35.0\% to 45.0\% &
The generalized Colab adapter/evaluation path runs end to end and can change
measured outputs. \\
\bottomrule
\end{tabularx}
\end{table}

\resultbox{\textbf{Fine-tuning conclusion.} A narrow adapter can shift Qwen's
Fashion-MNIST behavior enough to improve the paired result.  The new shirt
bias is not a footnote; it is exactly the kind of error redistribution that
has to be checked before calling a specialist model successful.}

\chapter{Limitations and Next Experiments}

\section{Limits of the present evidence}
The comparison is useful because it covers the full represented
model--benchmark rectangle, but it is still a study of the artifacts available
here.  Some runs use modest sample sizes.  Some rows have many ties.  One
detection row is zero for every model.  Normalization helps combine unlike
metrics, but it cannot turn a small or non-discriminating task into strong
evidence.

The telemetry has a different limit.  It honestly records what this harness
observed, but the execution paths mix hosted API calls and local model runs.
That can guide this repository's workflow, but it cannot establish intrinsic
inference efficiency across deployment modes.

The fine-tuning experiment is narrower again.  It evaluates one adapter family
on one classification task and one held-out sample.  The accuracy gain is worth
reporting, but the class-level regressions are exactly why more repetitions are
needed.

\section{Next experiments worth running}
\begin{enumerate}[leftmargin=22pt]
\item Repeat the balanced Fashion-MNIST QLoRA training with multiple seeds and
larger held-out evaluations; report mean accuracy and confusion matrices.
\item Diagnose the all-zero \code{flickr30k\_entities} detection result by
checking prompt format, coordinate normalization, parsing, and example-level
outputs before using it in model selection.
\item Fine-tune Qwen independently on one captioning task and one QA task.
This tests whether its observed captioning strength can be enhanced without
moving a narrow classifier's bias problem into a new setting.
\item Record deployment mode and repeated latency trials explicitly, so the
quality--speed comparison distinguishes API round-trip latency from local GPU
generation.
\item Increase benchmark sample counts on rows with very large observed
spread, then regenerate the bootstrap and pairwise analyses.
\end{enumerate}

\clearpage
\chapter*{Reproduction Notes and Sources}
\addcontentsline{toc}{chapter}{Reproduction Notes and Sources}

\section*{Reproducing the comparison}
\addcontentsline{toc}{section}{Reproducing the comparison}
The main tables and graphs in this report are generated from the top-level
JSON reports under \code{results/}.  From the repository root, run:

\begin{verbatim}
python comparison\build_report.py --results-dir results \
  --output-dir comparison\output --bootstrap-samples 5000
\end{verbatim}

This produces the CSV tables and plots used in Chapters 4--7.  The nested
fine-tuning reports are intentionally not included by that command.  They are
examined separately so the baseline comparison stays separate from the
task-adapted model.

\section*{Local sources used}
\addcontentsline{toc}{section}{Local sources used}
\begin{longtable}{>{\raggedright\arraybackslash}p{67mm}p{78mm}}
\toprule
\textbf{File or directory} & \textbf{Role in this report} \\
\midrule
\endhead
\path{comparison/output/model_summary.csv} &
Aggregate normalized scores, raw means, ranks, and first-place totals. \\
\path{comparison/output/family_summary.csv} &
Family-specific leader comparisons. \\
\path{comparison/output/pairwise_wins.csv} &
Head-to-head benchmark outcomes. \\
\path{comparison/output/bootstrap_summary.csv} &
Benchmark-resampled uncertainty intervals. \\
\path{comparison/output/telemetry_model_summary.csv} &
Timing, memory, and measured-coverage summary. \\
\path{models/_openai_vision.py} and \path{models/_hf_model.py} &
API and generic local inference pathways. \\
\path{models/_qwen2_5_vl.py}, \path{models/llava15_7b.py},
\path{models/paligemma_3b_mix_224.py} &
Family-specific wrapper behavior discussed in Chapter 3. \\
\path{benchmarks/_base_benchmark.py} &
Prompt, image, label, and evaluation contract. \\
\path{fine_tuning/train_qwen25vl_benchmark.py} &
QLoRA configuration and manifest-driven training behavior. \\
\path{results/fine-tuning/*.json} &
Stored 50-example Fashion-MNIST paired evaluations. \\
\path{fine_tuning_smoke_artifacts (1)/*.json} &
Downloaded 20-example Colab smoke comparison. \\
\bottomrule
\end{longtable}

\section*{External architecture references}
\addcontentsline{toc}{section}{External architecture references}
The architecture descriptions in Chapter 2 also use public model documentation
for the local open model families.  The API-served models are still described
only at the request/response boundary visible in this repository.

\begin{longtable}{>{\raggedright\arraybackslash}p{48mm}p{97mm}}
\toprule
\textbf{Reference} & \textbf{Use in this report} \\
\midrule
\endhead
\path{https://huggingface.co/papers/2502.13923} &
Qwen2.5-VL technical report summary: dynamic-resolution vision transformer,
window attention, and multimodal positional handling. \\
\path{https://huggingface.co/docs/transformers/v4.51.3/en/model_doc/qwen2_5_vl} &
Transformers implementation documentation for Qwen2.5-VL conditional
generation and configuration. \\
\path{https://huggingface.co/docs/transformers/v4.48.1/en/model_doc/llava} &
Transformers LLaVA documentation: CLIP vision encoder, MLP projection, and
autoregressive language-model generation. \\
\path{https://ai.google.dev/gemma/docs/paligemma/model-card} &
Google PaliGemma model card: PaLI/Gemma lineage, SigLIP vision component, and
Gemma language model component. \\
\bottomrule
\end{longtable}

\section*{Closing observation}
\addcontentsline{toc}{section}{Closing observation}
The main lesson is not simply that one model wins, and it is not that rankings
are meaningless.  The evaluated models do have an aggregate order in this
result set, and \code{gpt-5.5} leads it.  But the useful engineering decision
usually happens one level lower: Qwen's captioning behavior, PaliGemma's QA
result, the latency frontier, and the adapter's shirt bias all change what
``best'' means for a concrete application.  The value of the shared harness is
that it makes those differences visible.

\end{document}
